\section{Related Work}
\label{sec:related_work}

\textbf{Multi-objective optimization in drug discovery.}
Multi-objective molecular optimization has been explored through evolutionary algorithms \citep{deb2002nsga2,zhang2007moead}, generative models \citep{jin2020multiobjective,xie2021mars}, and benchmark suites such as GuacaMol \citep{brown2019guacamol}.
These approaches typically optimize over a generative model's output space, proposing novel molecules that balance competing objectives.
Our setting differs: the candidate pool is fixed, and the task is ranking (prioritization) rather than generation.
Scalarization, Pareto, and lexicographic methods for multi-objective ranking have been studied in the broader optimization literature \citep{emmerich2018tutorial}, but their application to peptide candidate triage with multiple estimated endpoints is largely unexplored.

\textbf{Peptide property prediction.}
Computational tools for individual peptide properties are well-established: antimicrobial activity prediction \citep{veltri2018ampscanner}, toxicity prediction \citep{rathore2024toxinpred3}, and stability estimation \citep{sharma2014hlp}.
Protein language models such as ESM-2 \citep{lin2023esm2} provide general-purpose sequence embeddings that improve downstream predictors.
However, most tools predict a single endpoint; few address the joint ranking problem across multiple developability criteria.
Learning-to-rank methods \citep{liu2009learningtorank,burges2010ranknet} offer a principled framework for combining multiple signals into a single ranking, but their application to peptide triage is unexplored.

\textbf{Autonomous ML research.}
The autoresearch framework \citep{karpathy2026autoresearch} demonstrated that a language-model agent can iteratively improve a language model training script through autonomous code modification against a fixed evaluation harness.
Subsequent work has explored language-model-guided program search in broader contexts: EvoPrompting \citep{chen2023evoprompting} for neural architecture search, FunSearch \citep{romeraparedes2024funsearch} for mathematical discovery, and The AI Scientist \citep{lu2024aiscientist} for open-ended scientific investigation.
ML-Agent-Bench \citep{huang2024mlagentbench} provides a benchmark for evaluating language agents on ML experimentation tasks.
Our companion paper \citep{wijaya2026autonomousagentdiscoversmolecular} applied the autoresearch paradigm to molecular transformer architecture search across SMILES, protein, and NLP domains with 3{,}106 experiments, finding that the value of architecture search is domain-dependent.

This paper extends the autoresearch paradigm along two axes: from architecture search (single editable file, single metric) to ranking policy search (single editable file, multi-objective evaluation), and from language modeling to peptide candidate triage.
The one-editable-file discipline and fixed-harness evaluation are preserved; the evaluation structure changes from scalar val\_bpb to a three-oracle ensemble.
